\chapter{Analytical procedures} % (fold)
\label{cha:analytical_procedures}

\begin{quote}
`When a truth is necessary, the reason for it can be found by analysis, that is, by resolving it into simpler ideas and truths until the primary ones are reached.'\\
Gottfried Wilhelm von Leibniz (1646-1716)
\end{quote}

As we are about to enter the arena of \hlblue{substantive analytical procedures}\index{substantive analytical procedures} (SAP)\index{SAP}, it is wise to first review the professional literature on these procedures. 

The main source of information is the International Standard on Auditing 520, “Analytical Procedures,” and its US equivalent AU-C 520 bearing the same title. Similar to what we remarked when discussing Standard 530, there are few textual differences between the two. The ISA uses the word ‘shall’ where the AU says ‘should.’ More importantly, the AU includes a documentation requirement that is absent from the ISA.

However, these Standards are not elaborate. There are seven requirement paragraphs in the ISA and eight in the AU, followed by 21 Application and Guidance paragraphs in the ISA, and 30 in the AU. The AICPA issued two additional guides that contain interesting clarifications and examples: the AICPA Analytical Procedures Audit Guide ("AAG-ANP")\index{AAG-ANP} and the AICPA Guide to Data Analytics ("AAG-DATA")\index{AAG-DATA}. In this chapter, we skim through the contents of this professional literature, particularly those parts that describe analytical procedures aimed at obtaining substantive audit evidence, the so-called SAPs. We are particularly interested in the extent to which they support the use of advanced statistical procedures such as regression analysis. Recall from Chapter \ref{cha:sampling} that the review of the sampling standard was cumbersome because it supported both statistical and non-statistical techniques. We note a similar problem with SAPs.

The scope of the standard is SAPs, as well as the analytical procedures that auditors are required to carry out as an overall review of the financial statements at the end of their audit. In this chapter, we limit our review to the requirements and guidance for SAPs.

Analytical procedures can also be used for risk-assessment. If this is the case, the requirements of Standard 315 prevail. It may be the same kind of procedure that an auditor employs to obtain substantive audit evidence; the only difference is that a risk assessment procedure mainly focuses on the identification of unusual patterns and outliers to identify risk areas and refine the audit program. The question to what extent the auditor has gained comfort when these procedures yield nothing unusual or remarkable is not answered; it is what risk assessment procedures distinguish from SAPs. Standard 330 includes the requirements and guidance for assessment. The audit procedures that the auditor then chooses to carry out in response to these risks may include the same advanced statistical procedures that would be employed as SAPs.

\bigskip
Now what are SAPs? The definition in the Standards\footnote{ISA 520.4, AU-C 520.4} emphasizes three important concepts. First, it is the \hlblue{evaluation of financial information}, which suggests that analytical procedures will be used to understand or test financial statement relationships or balances.

Secondly, the evaluation takes place through the analysis of \hlblue{plausible relationships}. 
This means that the auditor is not permitted to incorporate any data in the evaluation, even if it shows a significant correlation with the financial data under examination. 
A logical justification for the relationship between these two aspects must be established. This also implies that the auditor cannot have an algorithm that selects the data points to be used in the analysis; the auditor must be capable of justifying the choice to incorporate specific data in the analysis.

Third and final, the standards suggest that \hlblue{both financial and non-financial data} can be used to form the relationships of financial information, and therefore, also in forming expectations. 
For example, in the analysis of interest expenses, the auditor may use the financial account of long-term debt, which is financial, whereas in the analysis of payroll expenses, the auditor may use the personnel head count, which is a nonfinancial variable.

\bigskip
The definition of analytical procedures includes 'the investigation of identified fluctuations or relationships that are inconsistent with other relevant information, or that differ from expected values by a significant amount'.\footnote{ISA / AU-C 520.7} This implies an understanding of what can reasonably be expected and involves a comparison of the recorded book values with the auditor’s expectations and an understanding of these differences.

Paragraph 520.A1 provides some examples of comparisons: auditors may compare an account balance with that of prior periods, budgets and forecasts, and industry information. Paragraph 520.A2 states that the auditor expects information to conform to a predictable pattern. Paragraph 520.A3 mentions the entire range, from a simple comparison to complex analyses using advanced statistical techniques applied at all aggregation levels. Later, we specifically mention trend and ratio analyses, which are techniques that can be carried out statistically and non-statistically.






\section*{Learning objectives}
At the end of this chapter you will be able to

\begin{compactitem}

% - **Remember**:
\item define\marginnote{remember} a substantive analytical procedure (SAP) as outlined in auditing standards.
\item identify the different phases in planning an SAP.
\item list the steps to obtain relevant data while considering reliability, different sources, and levels of disaggregation.

% - **Understand**:
\item explain\marginnote{understand} how effectiveness risk and efficiency risk influence the determination of the acceptable difference in SAPs.
\item describe the complexity of model development and the various choices involved in predictive modeling.
\item understand the importance of documentation throughout all phases of an SAP.

% - **Apply**:
\item use\marginnote{apply} reliability and disaggregation principles to obtain relevant data for SAP.
\item evaluate the precision of an expectation by comparing it to the recorded amount or ratio.
\item take appropriate actions to increase precision and identify possible and probable causes of deviations.

% - **Analyze**:
\item assess\marginnote{analyze} the role of effectiveness risk and efficiency risk in SAP outcomes.
\item break down the components of predictive modeling to understand their influence on the precision of the expectation.

% - **Evaluate**:
\item judge\marginnote{evaluate} whether the objectives of the SAP have been achieved based on the analysis of deviations.
\item recommend improvements to the SAP process for enhancing precision and achieving audit goals.

% - **Create**:
\item develop\marginnote{create} a comprehensive SAP plan, incorporating the phases of planning, data collection, modeling, and evaluation.
\item design strategies for documenting the SAP process to meet professional auditing standards.
\end{compactitem}

\begin{mdframed}[hidealllines=true,backgroundcolor=KPMG_blue!20]
\section*{Case: On the Go Stores}

On the Go Stores have 23 convenience stores located in the Southeast. Included in the 23 stores are 5 new stores (Nos. 1, 4, 10, 13, and 22) that opened during the year. Operations vary by geographic location and the mix of products sold.

The location of a store is based on several factors, such as competition and the economic environment of the location. Store Nos. 2, 4, 6, 8, 9, 11, 13, 15, 17, 18, 20, 21, and 23 are considered to be in favorable locations.

Typically, a store's operations do not change much unless a new product line is introduced, such as selling gas, offering check-cashing services, or selling lottery tickets. The mix of products and services can vary, and the most important factor is whether the store sells gasoline. (Store Nos. 5, 6, 7, 8, 14, 15, 16, 17, 18, 19, 20, and 21 sell gasoline.) These additional product lines typically affect the volume of customers as well as the number of full-time employees.

There have been no significant event or accounting changes, except for the opening of new stores. Industry and economic factors along with management incentives remain the same compared to prior years. On the Go Stores is a repeat audit engagement. Materiality is 150,000.

\end{mdframed}

\begin{mdframed}[hidealllines=true,backgroundcolor=KPMG_blue!20]
\section*{Case: US SteamCo}
\label{case:ussteamco_ch06}

US SteamCo is a public utility in the northeastern United States. It produces steam and pumps it under high pressure to apartment buildings to provide heat in winter and power for air conditioners in summer. US SteamCo's regulator, the local public service commission, determines pricing, which varies by season and demand, and includes various fees, some fixed, and adjustments for variations in US SteamCo's fuel costs. The schedule is different for November to May (heating season) and for June to October (cooling season). Pricing schedules are revised at least annually.

You have performed other audit procedures to provide assurance regarding revenues, including the following:

\begin{compactitem}
\item Tests of the effective operation of controls over revenues and revenue adjustments (all relevant assertions).
\item Confirmation of accounts receivable from customers (occurrence, accuracy, cutoff).
\item Various procedures to audit cash receipts regarding amounts billed (occurrence).
\end{compactitem}

\end{mdframed}

\section*{Objective}
Apply the requirements and guidance from professional standards when performing SAPs in response to the relevant \hlblue{risk of material misstatement}\index{risk of material misstatement} (\hlblue{RMM}\index{RMM}) for the current year's sales. 

\section{Steps in performing a substantive analytical procedure} % (fold)
\label{sec:steps_in_performing_a_substantive_analytical_procedure}

The execution of SAPs involves many steps, and is broadly divided into four phases. The process is not strictly linear; some steps may be executed simultaneously, and sometimes it is necessary to revise earlier steps based on the outcome of later steps. 

\bigskip
The phases and steps are as follows:

\begin{enumerate}
\item Plan the SAP.
\begin{compactitem}
\item Determine the financial statement item or account and related assertions for which the SAP is to provide audit evidence and the specific objectives of the SAP.
\marginnote{Section \ref{sub:account_and_assertions}}
\item Identify the assessed risks of material misstatement to which the SAP is intended to respond (desired level of assurance).
\marginnote{Section \ref{sub:assessed_risk_of_material_misstatement}}
\item Determine whether a SAP is suitable to address the RMM.
\marginnote{Section \ref{sub:suitability_of_the_sap}}
\item Determine the plausibility of the relationship.
\marginnote{Section \ref{sub:plausibility_of_the_relationship}}
\item Determine the desired precision of the expectation.
\marginnote{Section \ref{sub:desired_precision}}
\item Determine the SAP (for example, trend analysis, ratio analysis, non-statistical predictive modelling, regression analysis) that is most likely to meet the objectives.
\marginnote{Section \ref{sub:choice_of_procedure}}
\end{compactitem}
\item Obtain the data from which the expectation of recorded amounts or ratios is to be developed, including
\begin{compactitem}
\item the relevance of the data;
\marginnote{Section \ref{sub:relevance_of_the_data}}
\item the reliability of the data;
\marginnote{Section \ref{sub:reliability}}
\item sources of the data for those variables; and
\marginnote{Section \ref{sub:source_of_the_data}}
\item levels of disaggregation of the data.
\marginnote{Section \ref{sub:disaggregation}}
\end{compactitem}
\item Perform the SAP.
\begin{compactitem}
\item Develop the model to be used, when applicable.
\marginnote{Section \ref{sub:model_development}
% , and \ref{sec:modeling} to \ref{sec:testing_significance}
}
\item Evaluate whether the expectation is sufficiently precise and, if not, the actions to take to increase the precision.
\marginnote{Section \ref{sub:evaluation_of_precision}}
\item Determine the amount of difference from the expectation that can be accepted without further investigation.
\marginnote{Section \ref{sub:acceptable_difference}}
\item Develop the expectation of the recorded amount or ratio.
\marginnote{Section \ref{sub:expectation_development}}
\item Compare the expectation to the recorded amount or ratio.
\end{compactitem}
\item Evaluate and respond to the results of the SAP.
\marginnote{Section \ref{sec:evaluation_and_response}}
\begin{compactitem}
\item Determine whether the difference between the recorded amount and our expectation is significant.
\item Investigate any significant difference, identifying possible and probable causes.
\item Determine whether the SAP has identified a misstatement and evaluate any misstatement.
\item Conclude on whether our objectives of the SAP have been achieved. If the objectives have not been achieved, plan and perform different procedures to achieve those objectives.
\end{compactitem}
\end{enumerate}

The final section of this chapter, Section \ref{sec:documentation}, focuses on documentation. It is important to emphasize that documentation plays a crucial role throughout all preceding phases.

\section{Planning the SAP} % (fold)
\label{sec:planning_the_sap}

\subsection{Account and assertions} % (fold)
\label{sub:account_and_assertions}

We start with determining the financial statement item or account and related assertions for which the SAP is to provide audit evidence and the specific objectives of the SAP. 

The \textit{Case: US SteamCo}, is intended to be a source of audit evidence regarding revenue earned by the company by providing steam for heating and cooling.\footnote{AAG-DATA B.16} The assertions addressed by this procedure are as follows:   
\begin{compactitem}
\item \textit{Occurrence}. All revenue transactions that have been recorded have occurred and pertain to the company.
\item \textit{Completeness}. All revenue transactions and events that should have been recorded have been recorded.
\item \textit{Accuracy}. Amounts and other data related to revenue transactions and events have been recorded properly.
\item \textit{Cutoff}. Transactions and events related to revenue have been recorded in the correct accounting period.\footnote{"Events" refers to, for example, adjustments to revenues by means of journal entries}
\end{compactitem}

The objective of the \textit{Case: On the Go Stores} is to examine sales (including merchandise sales and gas sales) analytically to determine the potential for overstatement and to address the auditor's objectives for testing occurrence and existence.\footnote{AAG-ANP A.52}


% subsection account_and_assertions (end)

\subsection{Assessed risk of material misstatement} % (fold)
\label{sub:assessed_risk_of_material_misstatement}

We then identify the assessed \hlblue{risk of material misstatement}\index{risk of material misstatement} (\hlblue{RMM})\index{RMM} to which the SAP is intended to respond. 
The RMM is 'the risk that the financial statements are materially misstated prior to the audit. The risk is a function of the susceptibility of the relevant assertion to misstatement and the effect of a misstatement on the financial statements.'\footnote{ISA / AU-C 315} The RMM refers to the likelihood that the financial statements of an entity contain errors or omissions that could affect the decisions of users of those statements. It includes both inherent risk and control risk, which are the risks associated with the nature of the business and the effectiveness of internal controls, respectively.

The RMM determines the desired level of assurance. For example, when the inherent risk is significant, and we place no reliance on controls, the desired level of assurance is 95\%. Auditing firms set lower levels of assurance when inherent risk is lower (\emph{elevated} or \emph{base} rather than \emph{significant}), and/or when the auditor places reliance on controls.

For the \textit{Case: US SteamCo}, we perform procedures appropriate in the circumstances encountered in this audit to respond to the presumption that risks of fraud exist in revenue recognition.\footnote{AAG DATA B.17}

This analytical procedure is designed to respond to a moderate level of risk of material misstatement of revenue.\footnote{AAG DATA B.18} 
In addition to performing the SAP, we also plan the following procedures:

\begin{compactitem}
\item Tests of the effective operation of controls over revenues and revenue adjustments (all relevant assertions)
\item Confirmation of accounts receivable from customers (occurrence, accuracy, cutoff)
\item Various procedures to audit cash receipts regarding amounts billed (occurrence) 
\end{compactitem}

The extent of the substantive procedures (confirmation and cash proof) depends on the level of assurance obtained from the SAP.

% subsection assessed_risks_of_material_misstatement (end)

\subsection{Suitability of the SAP} % (fold)
\label{sub:suitability_of_the_sap}

Suitability is a judgement call. As auditors begin utilizing SAPs, they develop an understanding of effective and ineffective approaches. The guidance in paragraph 520.A6\footnote{ISA 520.A6, AU-C 520.A6} provides some general suggestions. For example, accounts or classes of transactions consisting of large volumes of transactions that tend to be predictable over time lend themselves well to SAPs. Think of monthly sales, or the relationship between monthly revenue and monthly cost of goods sold, particularly when there are good cutoff procedures at the end of each month. 

The second part of paragraph 520.A6 is important. It states that 'the auditor should have the expectation that relationships among the data exist and continue in the absence of known conditions to the contrary.' Let us clarify this using an example. Assume that we have observed a steady relationship between revenue and the cost of goods sold over the last few years. In the year under audit, we expect this relationship to continue, and we have no knowledge of current or recent events that defeat this expectation. In this case, it is clear that the SAP is suitable. But what if our client starts consolidating a new entity that has markedly different profit margins? The expectations based on prior years’ relationships may be off by a significant amount. Based on our business understanding, we should not proceed with this analysis, as it is likely that some of our expectations may be off by a significant amount. For improved results, consider splitting the data into segments and conducting distinct analyses for each segment.

Paragraph 520.A6 also suggests that the suitability depends on the effectiveness of the procedure. This means that we should only consider an SAP if we expect it to be sufficiently precise to provide audit evidence. We return to this issue when discussing the precision of the expectation.\footnote{Section \ref{sub:evaluation_of_precision}}

According to paragraph 520.A7, in some cases even an unsophisticated predictive model may be effective as an analytical procedure. Consider, for example, the month-to-month comparison of an expense category, such as office rent. Typically, this is not an expense category that should display large fluctuations; it is more likely that the amount is constant over time, apart from an agreed-upon increase on a certain date. Our expectation is that the month-over-month differences are zero, and if they are not, we would examine it. This is also an excellent example of an analytical procedure that works on a small volume of transactions.

% subsection suitability_of_the_sap (end)

\subsection{Plausibility of the relationship} % (fold)
\label{sub:plausibility_of_the_relationship}

Analytical procedures operate on the premise that plausible relationships among data are expected to exist. It is crucial to assess the factors contributing to plausibility, as certain data types may appear related when, in reality, they are not. Such misperceptions can lead auditors to draw incorrect conclusions. On the other hand, unexpected relationships can offer valuable evidence when subjected to careful scrutiny.

Practical examples of plausible relationships are the following.

\begin{compactitem}
\item \textit{Inventory Turnover}: In the retail industry, the relationship between sales revenue and inventory levels is usually plausible. Higher sales should result in lower inventory levels as products are sold. However, if the analytical procedure reveals an unexpected relationship of increasing inventory levels despite rising sales, this could be an indication of potential overstocking, slow-moving items, or other inventory management issues that require further investigation.
\item \textit{Fixed Assets and Depreciation}: For a manufacturing company, there should be a plausible relationship between the cost of fixed assets and the corresponding depreciation expense. The analytical procedure could involve comparing the historical cost of fixed assets with the depreciation expense recognized over time. If the relationship appears unusual, such as a sudden spike in depreciation relative to the cost of assets, it may raise concerns about potential impairment or irregular accounting practices.
\item \textit{Accounts Payable and Purchases}: In most cases, the accounts payable balance is related to the level of purchases made by the company. Higher purchases should result in a corresponding increase in accounts payable. However, if the analytical procedure indicates a significant disconnect between the two, it could indicate potential discrepancies in recording payables or unrecorded liabilities.
\item \textit{Revenue and Cash Collections}: For a service-based company, there should be a plausible relationship between reported revenue and cash collections from customers. A notable disparity between revenue and cash collected may signal issues with revenue recognition practices or the effectiveness of the company's collections process.
\item \textit{Payroll Expenses and Employee Count}: In businesses with relatively stable employee compensation structures, there should be a plausible relationship between payroll expenses and the number of employees. If payroll expenses appear unrelated to the workforce size, it may indicate potential issues with unauthorized or inaccurate payroll entries.
\end{compactitem}

\bigskip
In each of these examples, we use analytical procedures to evaluate the plausibility of relationships between key financial data. Unexpected or implausible relationships can raise red flags, prompting us to dig deeper into the underlying transactions and accounts to understand the reasons behind these anomalies. This scrutiny helps identify potential risks of material misstatement and ensures a more comprehensive and effective audit process.

In \textit{Case: On the Go Stores} the auditor has selected a number of relevant independent variables, that is, those factors that the auditor knows from experience with the client and industry will be useful predictors of sales at each store. These independent variables are the level of inventory (merschandise plus gas) at the store, the number of staff at the store, whether the store opened or closed during the year, whether the store sells gas, and square feet of floor space at each store.

% subsection plausibility_of_the_relationship (end)

\subsection{Desired precision} % (fold)
\label{sub:desired_precision}

Precision refers to how closely the auditor's expectation aligns with the actual correct amount. The level of precision obtained from analytical procedures can vary based on the audit stage or the purpose of the procedure. When used as substantive tests, analytical procedures demand higher precision than those employed during the planning phase, if the analytical procedure is the only substantive procedure performed to address the RMM. The effectiveness of these procedures relies on their precision and intended purpose, and several factors influence their level of precision.

Professional literature mentions for different types of analysis, with increasing inherent precision:
\begin{compactitem}
	\item \textit{Trend analysis.} The analysis of changes in an account balance over time.\footnote{AAG-ANP 1.26}
	\item \textit{Ratio analysis.} The comparison of relationships between financial statement accounts over time.\footnote{AAG-ANP 1.30}
	\item \textit{Reasonablesness testing.} The analysis of account balances or changes in account balances within an accounting period that involves the development of an expectation based on financial data, nonfinancial data, or both.\footnote{AAG-ANP 1.33}
	\item \textit{Regression analysis.} This is the use of statistical models to quantify the auditor's expectation in monetary terms, with measurable risk and precision levels.\footnote{AAG-ANP 1.35}
\end{compactitem}

\bigskip
We argue in Section \ref{sub:acceptable_difference} that the acceptable difference is driven by the precision of the expectation and that this precision is difficult to control. Rather than specifying desired precision upfront, we measure the level of assurance obtained from the SAP to determine the extent of other substantive procedures.

% subsection desired_precision (end)

\subsection{Choice of procedure} % (fold)
\label{sub:choice_of_procedure}

When choosing between the different types of SAPs mentioned in the previous section, we should consider several factors to ensure the selected procedure is appropriate and effective. Some key factors to consider are:

\begin{compactitem}
\item \textit{Nature of the account or assertion}: The characteristics of the account or assertion being tested will guide our choice. Certain accounts may be better suited for specific analytical procedures. For example, a trend analysis may be more appropriate for revenue recognition, while a ratio analysis could be relevant for analyzing liquidity.

\item \textit{Data availability and reliability}: The availability and reliability of data play a crucial role in the selection process. If there is insufficient or questionable data, certain procedures may not be viable options. We should verify the source and quality of the data before choosing an analytical procedure.

\item \textit{Complexity of operations}: The complexity of the entity's operations may influence the choice of procedure. More complex operations might require sophisticated analytical techniques like regression analysis to capture intricate relationships between variables.

\item \textit{Materiality and risk assessment}: The materiality of the account and the assessed level of risk are important considerations. High-risk areas or material account balances may require more robust procedures, such as regression analysis, to provide sufficient assurance.

\item \textit{Trending data}: Trend analysis is particularly useful when comparing financial data over time to identify significant fluctuations or patterns. It can be beneficial in assessing the entity's financial performance and identifying potential risks.

\item \textit{Relationships between variables}: When there are clear relationships between different variables, ratio or regression analysis can be valuable in understanding how changes in one variable affect others. It helps identify potential inconsistencies or anomalies.

\item \textit{Objective and subjective factors}: Reasonableness tests are based on subjective judgments and the auditor's professional judgment. They are often used when there is no clear quantitative relationship between variables but where the auditor can assess the reasonableness of the figures based on experience and industry knowledge.

\item \textit{Time and resources}: The time and resources available for conducting the analytical procedure will impact the choice. Regression analysis, for example, can be time-consuming and require specialized knowledge and software.

\item \textit{Prior experience and success}: Prior experience with specific ana\-lyt\-i\-cal procedures and their effectiveness in identifying material mis\-state\-ments may influence the choice. Procedures that have been successful in previous audits may be preferred.

\item \textit{Regulatory and reporting requirements}: The applicable regulatory requirements and reporting standards may prescribe or recommend specific analytical procedures, and auditors should comply with these guidelines.

\end{compactitem}

\bigskip
Ultimately, we should consider a combination of these factors and exercise professional judgment to select the most suitable analytical procedure for the audit. It's essential to tailor our choice to the specific risks, complexities, and characteristics of the audit engagement to obtain sufficient and appropriate audit evidence for the financial statements.

% subsection choice_of_procedure (end)

% section planning_the_sap (end)

\section{Obtaining the data} % (fold)
\label{sec:relevance_and_reliability_of_the_data}

\subsection{Relevance of the data} % (fold)
\label{sub:relevance_of_the_data}

Considering the relevance and reliability of the information used in SAPs is crucial for the overall success and quality of the audit. It ensures that our conclusions are well-supported, risk assessments are accurate, and the audit provides valuable insights to stakeholders. By following the guidance in Standard 500 (the standard that deals with audit evidence and provides guidance on the verification of information provided by the entity), we can conduct thorough and effective analytical procedures, contributing to a robust and credible audit process.

Relevance relates to the logical connection with, or bearing upon, the purpose of the SAP and the risk of material misstatement and assertions under consideration. We evaluate the nature of the internal and external input data to determine whether they are sufficiently comparable and therefore relevant to the expectation developed for the assertions addressed by the SAP.

When inputs from several historical periods are used in the SAP, we consider the relevance of these specific periods. For example, if we use the trailing three fiscal periods of financial information, why? Why not two? Or five?

% subsection relevance_of_the_data (end)

\subsection{Reliability} % (fold)
\label{sub:reliability}

The second condition for a good SAP is the use of reliable data. A well-known phrase is: ‘garbage in -- garbage out.’ If the data used in the analysis were unreliable, the procedure would be less effective. The standard requires us to consider the source, comparability, nature, and relevance of available information. When considering the source, we should distinguish between internal and external data. We can assess the quality of internal data, for example, because we can investigate controls over the preparation of the data or whether the internal data have been subjected to audit testing. In addition, we may need to consider whether the internal data are independent of the account values that are subject to testing. If we cannot rely on controls, we may need to follow the provisions of Standard 500. Paragraph 500.7 establishes the requirements for verifying the reliability of the data. 

This, of course, is different from the external data. We may consider the reliability of the source; for example, government data from a government website or from a data provider usually have a different quality than social media data scraped from the internet.

Contamination of data affects the precision of expectations, for example, when the quality of cutoff procedures at the month end is different from that at the quarter- or year-end. As such, when using advanced statistical analysis techniques, the data quality is directly reflected in the precision of the model. This may imply that we do not need to consider reliability when statistical methods are used because noise in the input data translates directly into precision and, as a result, into the level of assurance obtained. Of course, if the model is less precise than desired, we can further examine the reliability of the input data. The Standards do not anticipate this, and there are no exceptions to statistical methods when it comes to assessing reliability. For nonstatistical procedures, caution should be exercised when assessing the effectiveness of the procedure. 

We consider the accuracy and completeness of data when evaluating its reliability, taking into account, for example, the objectives of the procedures for which the data is being used. Performing the following steps in sequence can aid in achieving the objective of the procedure:

\begin{compactitem}
\item Evaluate the data likely to be most relevant for effective and efficient performance of SAP.
\item Ascertain whether all the considered relevant data is readily available. If not, we explore viable steps to access that data. It may result in having only some of the most relevant data, in which case we reconsider if there is sufficient data with significant relevance to provide useful audit evidence.
\item When data of sufficient relevance is available, we make a preliminary assessment of its likely reliability for the purposes of the SAP. 
\end{compactitem}

% subsection reliability (end)

\subsection{Sources of the data} % (fold)
\label{sub:source_of_the_data}

We have the flexibility to choose from several data sources to create expectations. These sources include prior-period data (adjusted for expected changes), management's budgets or forecasts, industry data, and nonfinancial data. The selection of the data source directly affects the precision of our expectations for an account balance or class of transactions. Therefore, it is crucial to consider the source of information carefully when developing expectations, as it plays a significant role in achieving the desired level of assurance from the analytical procedure.

% subsection source_of_the_data (end)

\subsection{Disaggregation} % (fold)
\label{sub:disaggregation}

Data \hlblue{aggregation}\index{aggregation} and \hlblue{disaggregation}\index{disaggregation} refers to how account balances are combined or segmented for testing purposes, such as using daily, weekly or monthly rather than quarterly or annual data. 
Generally, the greater the level of disaggregation in the data used to form the expectation, the more precise the expectation becomes. 
This heightened precision leads to a higher level of assurance in detecting material misstatements. 
Disaggregation becomes especially crucial in complex or diversified business operations.

However, we must also evaluate the reliability of disaggregated data. For instance, more fine-grained data might be less reliable than quarterly or annual data due to lack of auditing or weaker controls. 
Therefore, we should exercise judgment in determining which precision factor carries more significance given the specific circumstances at hand.

% subsection disaggregation (end)

% section relevance_and_reliability_of_the_data (end)

\section{Performing the SAP} % (fold)
\label{sec:application_of_the_sap}

In the Perform phase we develop a model, evaluate its precision, determine the acceptable difference, and develop expectations.

The first two steps are executed on the \hlblue{training set}\index{training set}\index{set!training}, which is the dataset or sample used to construct the model. 
The model is then applied to the \hlblue{hold-out set}\index{hold-out set}\index{set!hold-out} to validate the theory, determine the acceptable difference and develop expectations. 
The division between the two sets is obvious in time-series analysis: the training set consists of data from prior periods, and the hold-out set contains the values for the current year under audit. 
When applying cross-sectional analysis, the necessity of dividing the two sets should not be overlooked in order to avoid circular reasoning. 
Because it is advantageous to maximize the number of observations in the training set, we usually divide the observations into ten sets of equal size, where each set is evaluated on a model constructed with the data from the other nine sets.

\subsection{Model development} % (fold)
\label{sub:model_development}

The assertion that predictive modeling is an art, rather than a science, highlights the complexity and subjectivity involved in the process. While there are established principles and methodologies for developing predictive models, we have significant discretion in making various choices that can impact the model's outcomes. These choices often involve a combination of technical expertise, creativity, and judgment, leading to the characterization of predictive modeling as an art.

Some examples of the choices we face in predictive modeling are:

\begin{compactitem}
\item \textit{Feature selection}: We must decide which independent variables (features) to include in the model. Choosing relevant and meaningful features is essential for accurate expectations. The process may involve domain knowledge, data exploration, and statistical techniques like correlation analysis.
\item \textit{Data preprocessing}: Before building the model, we need to preprocess the data. This involves handling missing values, dealing with outliers, normalizing or scaling variables, and encoding categorical variables. The choice of data preprocessing techniques can significantly impact the model's performance.
\item \textit{Hyperparameter tuning}: Many models have hyperparameters that need to be set before training, for example, the degree of polynomial features in the linear regression model. Determining the optimal values for these hyperparameters often requires experimentation and validation. The choice of hyperparameters can significantly affect the model's predictive ability.
\item \textit{Validation and evaluation metrics}: We must choose appropriate evaluation metrics to assess the model's performance. Common metrics include precision and r-squared. The choice of metric depends on the nature of the problem and the business objectives.
\item \textit{Handling imbalanced data}: In some predictive modeling scenarios, the data may be imbalanced, with one class significantly outnumbering the others. We must decide on strategies to handle this imbalance to avoid bias in the model's expectations.
\item \textit{Overfitting and underfitting}: We need to strike a balance between model complexity and generalizability. Overfitting occurs when the model performs well on the training set but poorly on the hold-out set. Underfitting happens when the model is too simple to capture the underlying patterns. The choice of model complexity is critical to avoid these issues.
\item \textit{Interpretability vs. performance}: Some predictive models, like linear regression analysis, are highly interpretable, making them easier to explain to stakeholders. Other models, like neural networks, offer superior performance but are less interpretable. We must consider the trade-off between interpretability and performance based on the audience's needs.
\end{compactitem}

\bigskip
In summary, predictive modeling involves a series of choices and trade-offs that require technical expertise, creativity, and judgment. While there are scientific principles underpinning predictive modeling, our ability to make informed decisions and interpret the results is a matter of professional judgement. Striking the right balance in these choices is essential to develop reliable and effective predictive models that meet the specific requirements of the problem at hand.

% subsection model_development (end)

\subsection{Evaluation of precision} % (fold)
\label{sub:evaluation_of_precision}

The next requirement is the precision evaluation. \hlblue{Precision}\index{precision} is a measure of the closeness of expectations to correct amounts. 
This is not the same as the difference between the \emph{recorded} amount and the expectation. 
The correct amount is unknown; the recorded amount is known, but may be incorrect. 
The assumption is that, if the recorded amount is not misstated, on average, the expectation equals the correct amount, but it may be off by some amount of imprecision.

In the example of share capital, precision is zero: we expect this year's amount to be exactly the same as last year's. Precision is small when analyzing an account such as office rent: we can either look at the monthly charges, where precision is zero for the month without an increase, and is a margin around a certain expected percentage increase in the month when rent was increased. Alternatively, we can look at a comparison between this year’s charge and the previous year’s charge, which is the last year’s charge plus a certain increase (perhaps tied to an index), and precision would be some margin around the index.

\bigskip
For less precise analytical procedures, it is much more difficult to assess precision, because it is a function of three factors\footnote{AAG 2.03}:
\begin{compactitem}
\item the nature of the account or assertion; some accounts, such as fixed asset additions, are less predictable than others, such as revenue.
\item reliability and other data characteristics. If the data used for creating expectations are inaccurate, the resulting expectations will also be inaccurate.
\item the inherent precision of the expectation method.
\end{compactitem} 

\bigskip
This inherent precision may be due to several factors. There could be other sources of information that further help form an expectation, but these were not considered when forming the expectation. There may be other models that describe the relationship better; for example, instead of using a linear regression model, we could consider a nonlinear regression or a time-series model. Each of these modeling techniques has advantages and disadvantages, and a more complicated model does not always produce better results. Finally, the inherent precision can be reduced by adding more data (but unlike sampling, this is not always easy) or by disaggregating the data. Compared to stratification in sampling, this is only an effective technique if the strata are significantly different. 

\bigskip
When an account subject to testing is not misstated, the expectation is, on average, equal to the recorded amount. In other words, the difference between the recorded amount and the expectation is expected to be zero.

When an account balance is analyzed using a predictive analytical procedure, there is some inherent inaccuracy (imprecision) in forming expectations. 
Therefore, in most cases, the expectation will not be perfectly equal to the recorded amount, even when the account balance is not misstated. 
However, when the account balance is not misstated, it is expected that the majority of the differences observed between the expectation and the recorded amount are within a certain range around zero. 
This range is referred to as the \hlblue{prediction interval}\index{prediction interval}\index{interval!prediction}. 
The width of the interval reflects the expectation precision.

Figure \ref{fig:precision} shows the prediction intervals for a more precise expectation and for a less precise expectation. 
Both intervals are centered on zero and cover 99\%\footnote{The choice of this percentage is explained on page \pageref{mar:controlling_efficiency_risk}.} of the expected differences. There is a 0.5\% chance that a difference is observed that is less than $C_{.005}$, and a 0.5\% chance that a difference is observed that is greater than $C_{.995}$.

\begin{figure} 
\begin{center}
\begin{tikzpicture}

\draw [<->, >=stealth] (0, 0) -- (10, 0) node[below] {difference}; 

% \draw [Parenthesis-Parenthesis] (1, 0.5) node[above] {$U_{.05}$} -- (3, 0.5) node[above] {$U_{.95}$};
\draw [Parenthesis-Parenthesis] (3.8, 1) node[above] {$C_{.005}$} -- (6.2, 1) node[above] {$C_{.995}$};
\draw [Parenthesis-Parenthesis] (2, 0.5) node[above] {$C_{.005}$} -- (8, 0.5) node[above] {$C_{.995}$};
% \draw [Parenthesis-Parenthesis] (7, 0.5) node[above] {$O_{.05}$} -- (9, 0.5) node[above] {$O_{.95}$};

% \draw (2, -0.1) node[below] {$-PM$} -- (2, 0.1);
\draw (5,  0.4)                 -- (5, 0.6);
\draw (5,  0.9)                 -- (5, 1.1);
\draw (5, -0.1) node[below] {0} -- (5, 0.1);
% \draw (8, -0.1) node[below] { $PM$} -- (8, 0.1);

\end{tikzpicture}
\caption{Prediction intervals for more precise (above) and less precise (below) expectations. These intervals cover 99\% of the expected differences.}

\label{fig:precision}
\end{center}
\end{figure}

Figure \ref{fig:precision} shows that when an expectation is precise, the observed difference between the recorded amount and the expectation is expected to be within a narrow range around zero. The expectation is `close' to the recorded amount. When the expectation is less precise, the difference observed may be within a wider range.

\bigskip
The precision of expectations is directly linked to the level of assurance provided by the SAP. When expectations are more precise, the procedure is more effective in identifying potential misstatements, and the resulting level of assurance provided by the procedure is higher.

To evaluate whether the expectation is sufficiently precise when performing an SAP,\footnote{AU-C520.A21, ISA 520.A14} it is appropriate for the auditor to consider whether SAPs are the only substantive procedures planned to address a particular risk of misstatement at the relevant assertion level, or whether the risk will be addressed through a combination of SAPs and tests of details. A less precise expectation may be appropriate when evidence obtained from performing the SAP is combined with audit evidence from performing tests of details, or when the target level of assurance is low because of a lower assessment of inherent risk and/or reliance on controls. However, a more precise expectation is necessary when an SAP is the only procedure planned to address a particular risk of misstatement for a relevant assertion, or when addressing assertions with a higher inherent risk and/or when no reliance can be placed on controls.

In the discussion of the \text{Case: On the Go Stores}, reference is made to the $r^2$ (R-squared) obtained from the model, the $t$ statistic of each of the predictor variables, and the standard error of the estimate.\footnote{ANG-ANP A.58}. Each of these quantities is discussed in more detail in Chapter \ref{cha:introduction_to_regression_analysis}.

\bigskip
Expectation and its related precision are the cornerstones of an SAP. Precision also plays a key role in the development of the acceptable difference. This will be discussed in the next section.

% subsection evaluation_of_precision (end)

\subsection{Acceptable difference} % (fold)
\label{sub:acceptable_difference}

The next requirement of auditing standard 520\footnote{520.05d} is to set an \hlblue{acceptable difference}\index{acceptable difference}\index{difference!acceptable}. 
The acceptable difference is 'the amount of any difference of recorded amounts from expected values that is acceptable without further investigation'. 
The rationale for this requirement is straightforward. 
Once we have created an expectation, we compare it with the recorded amount and calculate the difference between the two. 
Because the expectation is somewhat imprecise, the expectation is often not exactly equal to the recorded amount. 
Therefore, differences found could be consistent with the imprecision of the model used to develop the expectation. 
A larger difference could be indicative of the following:

\begin{compactitem}
\item misspecification of the model;
\item model lacking important predictors;
\item the recorded amount being materially misstated; or
\item a combination of these factors. 
\end{compactitem}

\bigskip
Which size difference requires further investigation? The answer to this question hinges on the risk that the auditor wishes to control, and our argument is similar to that discussed in Chapter \ref{cha:hypothesis_testing} regarding hypothesis testing. 

We begin with Step 1 in the nine-step approach\index{nine-step approach} from Chapter \ref{cha:hypothesis_testing}: the research question. Section \ref{sub:account_and_assertions} provides details, for example, whether the auditor is concerned about overstatements, understatements, or both.
Step 2 involves the development of the hypotheses. In the case of an SAP, we define two sets of hypotheses: the first null hypothesis is that the current year’s recorded amount is materially overstated versus the alternative that it is not, and the second null hypothesis is that the current year’s recorded amount is materially understated versus the alternative that it is not. The auditor chooses whether the test is directed toward detecting overstatements, understatements, or both.

The statistical assumptions for Step 3 are discussed in more detail in the next chapter. The test statistic (Step 4) is the difference between the recorded amount and the expectation.\footnote{When calculating the residuals from a linear model, the convention is to subtract the fitted values from the observed values. Mathematically, it is represented as $d = y - \hat{y}$, where $y$ is the recorded amount and $\hat{y}$ is the expectation.} The distribution of the test statistic (Step 5) is discussed in the next chapter.

Step 6 involves setting the significance level. Since we are testing two sets of hypotheses, we set a significance level for each; they do not necessarily have to be equal to each other. For example, we may assess the risk of material overstatement at a different level from that of material understatement. As previously discussed, the combination of the risk of a material misstatement and the assessment of internal controls drives the required level of assurance, which is the complement of the significance level.

The determination of an acceptable difference is similar to the determination of a critical region (Step 7) in hypothesis testing. However, in the case of an SAP, we should decide which risk we want to control: the risk of concluding that a material misstatement exists when it does not, or the risk of concluding that a material misstatement does not exist when it does. 

We compute the value of the test statistic (Step 8), which is the difference between the recorded amount and the expectation in the case of an SAP. This step is discussed in more detail in Section \ref{sub:expectation_development}.
% 
Finally, the conclusion of the test (Step 9) is presented in Section \ref{sec:evaluation_and_response}.


\bigskip
We use performance materiality ($M$) in the design of a predictive method. When the account balance is overstated by $M$, we expect the difference between the recorded amount and the expectation to be in a certain range around $M$. A similar range can be constructed when the account is understated by $M$: in that case the differences are expected to be in a certain range around $-M$. Figure \ref{fig:expected_difference_precise_model} shows these respective ranges. 

If the account is not misstated, we expect most (99\%) differences to be between $C_{.005}$ and $C_{.995}$, whereas if the account is understated by $PM$, most (90\%\footnote{ A different percentage is used here intentionally, as will be explained on page \pageref{mar:controlling_effectiveness_risk}.}) differences will be between $U_{.05}$ and $U_{.95}$, or if the account is overstated by $PM$, most (90\%) differences will be between $O_{.05}$ and $O_{.95}$.

\begin{figure} 
\begin{center}
\begin{tikzpicture}

\draw [<->, >=stealth] (0, 0) -- (10, 0) node[below] {difference}; 

\draw [Parenthesis-Parenthesis] (1, 0.5) node[above] {$U_{.05}$} -- (3, 0.5) node[above] {$U_{.95}$};
\draw [Parenthesis-Parenthesis] (3.8, 1) node[above] {$C_{.005}$} -- (6.2, 1) node[above] {$C_{.995}$};
\draw [Parenthesis-Parenthesis] (7, 0.5) node[above] {$O_{.05}$} -- (9, 0.5) node[above] {$O_{.95}$};

\draw (2, -0.1) node[below] {$-PM$} -- (2, 0.1);
\draw (5, -0.1) node[below] { 0} -- (5, 0.1);
\draw (8, -0.1) node[below] { $PM$} -- (8, 0.1);

\draw (2,  0.4) -- (2, 0.6);
\draw (8,  0.4) -- (8, 0.6);
\draw (5,  0.9) -- (5, 1.1);

\end{tikzpicture}
\caption{Expected difference for precise model, when the account is understated by $PM$, correct, or overstated by $PM$}

\label{fig:expected_difference_precise_model}
\end{center}
\end{figure}

It is reasonable to assume that if a difference is observed between $C_{.005}$ and $C_{.995}$, the account is not materially misstated. In addition, when the difference observed is between $U_{.05}$ and $U_{.95}$, it is likely that the account is materially understated, and when the difference is between $O_{.05}$ and $O_{.95}$ it is likely that the account is materially overstated. In addition, if the difference is less than $U_{.05}$, it is likely that the account is understated by an amount that is greater than $PM$, and when the difference is greater than $O_{.95}$, it is likely that the account is overstated by more than $PM$. 

It follows that we can choose the boundary between `acceptable' and `not acceptable' at either $U_{.95}$ or $C_{.005}$ for understatements, and at either $C_{.995}$ and $O_{.05}$ for overstatements.
However, two further questions arise:

\begin{enumerate}
    \item What if the difference observed is in the `gap' between $U_{.95}$ and $C_{.005}$, or in the `gap' between $C_{.995}$ and $O_{.05}$?
    \item What if the model is less precise and the ranges overlap?
\end{enumerate}

Figure \ref{fig:expected_difference_less_precise_model} displays the second question. Because the model is now less precise than what was previously the case, the ranges of the expected differences are wider. Keeping the percentages the same, the ranges now overlap, and rather than having a situation in which there are gaps, there are now ambiguous outcomes.

\begin{figure} 
\begin{center}
\begin{tikzpicture}

\draw [<->, >=stealth] (0, 0) -- (10, 0) node[below] {difference}; 

\draw [Parenthesis-Parenthesis] (0, 0.5) node[above] {$U_{.05}$} -- (4, 0.5) node[above] {$U_{.95}$};
\draw [Parenthesis-Parenthesis] (2.8, 1) node[above] {$C_{.005}$} -- (7.2, 1) node[above] {$C_{.995}$};
\draw [Parenthesis-Parenthesis] (6, 0.5) node[above] {$O_{.05}$} -- (10, 0.5) node[above] {$O_{.95}$};

\draw (2, -0.1) node[below] {$-PM$} -- (2, 0.1);
\draw (5, -0.1) node[below] { 0} -- (5, 0.1);
\draw (8, -0.1) node[below] { $PM$} -- (8, 0.1);

\draw (2,  0.4) -- (2, 0.6);
\draw (8,  0.4) -- (8, 0.6);
\draw (5,  0.9) -- (5, 1.1);

\end{tikzpicture}
\caption{Expected difference for a less precise model, when the account is understated by $PM$, correct, or overstated by $PM$}

\label{fig:expected_difference_less_precise_model}
\end{center}
\end{figure}

Ambiguity arises when a difference is observed between $C_{.005}$ and $U_{.95}$ or between $O_{.05}$ and $C_{.995}$. The differences of this size are consistent with both the assumption that the account is not misstated and the assumption that it is materially misstated. To decide whether a difference is acceptable or needs further investigation, we need to first discuss which risks are involved. 

\paragraph{Managing risks} The intervals above are directly linked to the acceptable differences that can be determined, and with the two risks an auditor incurs when concluding under uncertainty:

\begin{compactitem}
\item \textit{Effectiveness risk}: the risk of deciding that a difference does not require further examination when the account under review is materially misstated, and 
\item \textit{Efficiency risk}: the risk of deciding that a difference requires further examination when the account under review is not materially misstated.
\end{compactitem}

\bigskip
The acceptable difference can be determined by controlling either the detection or the efficiency risk. Neither risk can be controlled simultaneously. We now examine the effects of the choice of risk to be controlled.

Figure \ref{fig:method_detection_precise}\label{mar:controlling_effectiveness_risk} shows how the acceptable difference is determined when the detection risk is set at 5\% for both overstatements and understatements. 

\begin{figure} 
\begin{center}
\begin{tikzpicture}

\draw [<->, >=stealth] (0, 0) -- (10, 0) node[below] {difference}; 

\draw [loosely dotted] (1, 1) node[above] {$U_{.05}$} -- (3, 1) node[above] {$U_{.95}$};
\draw [loosely dotted] (7, 1) node[above] {$O_{.05}$} -- (9, 1) node[above] {$O_{.95}$};

\draw [Parenthesis-] (3, 1.75) node[above] {$U_{.95}$} -- (6.5, 1.75) node[above] {acceptable for understatements} -- (10, 1.75) ;
\draw [dashed, Parenthesis-Parenthesis] (3.8, 0.5) node[above] {$C_{.005}$} -- (6.2, 0.5) node[above] {$C_{.995}$};
\draw [-Parenthesis] (0, 2.5) -- (3.5, 2.5) node[above] {acceptable for overstatements} -- (7, 2.5) node[above] {$O_{.05}$};

\draw (2, -0.1) node[below] {$-PM$} -- (2, 0.1);
\draw (5, -0.1) node[below] { 0} -- (5, 0.1);
\draw (8, -0.1) node[below] { $PM$} -- (8, 0.1);

\end{tikzpicture}
\caption{acceptable differences for a precise model, focusing on effec\-tiveness risk}

\label{fig:method_detection_precise}
\end{center}
\end{figure}

If the account is materially overstated by an amount of $PM$, the auditor expects differences within the range of $O_{.05}$ to $O_{.95}$. In this scenario, differences up to $O_{.05}$ are acceptable. As a result, the auditor incurs a 5\% risk of failing to investigate the difference when the account is materially overstated. This risk is related to the procedure's \emph{effectiveness}. This is synonymous with stating that the auditor has set a desired level of assurance of 95\% regarding a material overstatement. Note that this interval is one-sided. Differences greater than $O_{.95}$ are indicative of a misstatement that is larger than $PM$. Similarly, if the account is materially understated by $PM$, differences greater than $U_{.95}$ are acceptable. 

In this scenario based on a precise model, the efficiency risk is actually less than 1\%, because the interval of $O_{.05}$ to $O_{.95}$ is wider than that from $C_{.005}$ to $C_{.995}$.

\bigskip
Figure \ref{fig:method_detection_less_precise} shows a situation with poorer precision, similar to Figure \ref{fig:expected_difference_less_precise_model}. The detection risk was maintained at 5\%.

\begin{figure} 
\begin{center}
\begin{tikzpicture}

\draw [<->, >=stealth] (0, 0) -- (10, 0) node[below] {difference}; 


\draw [loosely dotted] (0, 1) node[above] {$U_{.05}$} -- (4, 1) node[above] {$U_{.95}$};
\draw [loosely dotted] (6, 1) node[above] {$O_{.05}$} -- (10, 1) node[above] {$O_{.95}$};


\draw [Parenthesis-] (4, 1.75) node[above] {$U_{.95}$} -- (7, 1.75) node[above] {acceptable for understatements} -- (10, 1.75) ;
\draw [dashed, Parenthesis-Parenthesis] (2.8, 0.5) node[above] {$C_{.005}$} -- (7.2, 0.5) node[above] {$C_{.995}$};
\draw [-Parenthesis] (0, 2.5) -- (3, 2.5) node[above] {acceptable for overstatements} -- (6, 2.5) node[above] {$O_{.05}$};

\draw (2, -0.1) node[below] {$-PM$} -- (2, 0.1);
\draw (5, -0.1) node[below] { 0} -- (5, 0.1);
\draw (8, -0.1) node[below] { $PM$} -- (8, 0.1);

\draw (2,  2.4) -- (2, 2.6);
\draw (5,  0.4) -- (5, 0.6);
\draw (8,  1.65) -- (8, 1.85);

\end{tikzpicture}
\caption{acceptable differences for a less precise model, focusing on effectiveness risk}

\label{fig:method_detection_less_precise}
\end{center}
\end{figure}
 
Because the precision of the expectation is less than that in Figure \ref{fig:method_detection_precise}, the ranges of the expected differences are now wider. As a result, the range of acceptable differences ($U_{.95}$ to $O_{.05}$) is narrower. While the detection risk is kept at the original 5\%, it follows that if the account balance is not misstated, the probability of observing a difference that requires further investigation is increased: the range of acceptable differences is narrower than the range of expected differences when the account is not misstated of $C_{.005}$ to $C_{.995}$. As precision deteriorates, the range of acceptable differences becomes narrower, even up to the point where points $U_{.95}$ and $O_{.05}$ coincide, in which case every difference observed requires further work.

To control\label{mar:controlling_efficiency_risk} for efficiency risk, the acceptable difference is determined by the interval $C_{.005}$ to $C_{.995}$. A difference between these two values is to be expected given the precision of the model and should therefore be no surprise. If indeed the account balance is not misstated, the auditor then incurs a 1\% risk of observing a difference that is a significant unexpected difference that requires further investigation. This risk is related to \emph{efficiency}. 

Figure \ref{fig:method_efficiency_precise} shows the acceptable difference for a precise model determined by controlling the efficiency risk at 1\%. The acceptable difference ranges from $C_{.005}$ to $C_{.995}$. The level of evidence obtained for overstatements is greater than 99\%, since (at this precision) $C_{.995}$ is less than $O_{.05}$. 
% The increase in the level of assurance obtained is indicated by the dotted line, that indicates that the boundary of the acceptable difference for overstatements is decreased from $O_{.05}$ to $C_{.995}$. 

Likewise, the level of evidence obtained for understatements is greater than 95\% because (at this precision) $C_{.005}$ is greater than $U_{.95}$. The boundary for the acceptable difference on understatements is moved up from $U_{.95}$ to $C_{.005}$
% , as indicated by the dotted line
.

\begin{figure} 
\begin{center}
\begin{tikzpicture}

\draw [<->, >=stealth] (0, 0) -- (10, 0) node[below] {difference}; 

\draw [loosely dotted] (1, 1) node[above] {$U_{.05}$} -- (3, 1) node[above] {$U_{.95}$};
\draw [loosely dotted] (7, 1) node[above] {$O_{.05}$} -- (9, 1) node[above] {$O_{.95}$};

\draw [Parenthesis-] (3.8, 1.75) node[above] {$C_{.005}$} -- (6.9, 1.75) node[above] {acceptable for understatements} -- (10, 1.75) ;
\draw [dashed, Parenthesis-Parenthesis] (3.8, 0.5) node[above] {$C_{.005}$} -- (6.2, 0.5) node[above] {$C_{.995}$};
\draw [-Parenthesis] (0, 2.5) -- (3.1, 2.5) node[above] {acceptable for overstatements} -- (6.2, 2.5) node[above] {$C_{.995}$};

\draw (2, -0.1) node[below] {$-PM$} -- (2, 0.1);
\draw (5, -0.1) node[below] { 0} -- (5, 0.1);
\draw (8, -0.1) node[below] { $PM$} -- (8, 0.1);

\end{tikzpicture}
\caption{acceptable differences for a precise model, focusing on efficiency risk}

\label{fig:method_efficiency_precise}
\end{center}
\end{figure}

If the expectation precision is poor, the acceptable difference determined by controlling the efficiency risk is wider. Figure \ref{fig:method_inefficiency_less_precise} illustrates this situation.

\begin{figure} 
\begin{center}
\begin{tikzpicture}

\draw [<->, >=stealth] (0, 0) -- (10, 0) node[below] {difference}; 

\draw [loosely dotted] (0, 1) node[above] {$U_{.05}$} -- (4, 1) node[above] {$U_{.95}$};
\draw [loosely dotted] (6, 1) node[above] {$O_{.05}$} -- (10, 1) node[above] {$O_{.95}$};


\draw [Parenthesis-] (2.8, 1.75) node[above] {$C_{.005}$} -- (6.4, 1.75) node[above] {acceptable for understatements} -- (10, 1.75) ;
\draw [dashed, Parenthesis-Parenthesis] (2.8, 0.5) node[above] {$C_{.005}$} -- (7.2, 0.5) node[above] {$C_{.995}$};
\draw [-Parenthesis] (0, 2.5) -- (3.6, 2.5) node[above] {acceptable for overstatements} -- (7.2, 2.5) node[above] {$C_{.995}$};

\draw (2, -0.1) node[below] {$-PM$} -- (2, 0.1);
\draw (5, -0.1) node[below] { 0} -- (5, 0.1);
\draw (8, -0.1) node[below] { $PM$} -- (8, 0.1);

\draw (2,  2.4) -- (2, 2.6);
\draw (5,  0.4) -- (5, 0.6);
\draw (8,  1.65) -- (8, 1.85);

\end{tikzpicture}
\caption{acceptable differences for a less precise model, focusing on efficiency risk}

\label{fig:method_inefficiency_less_precise}
\end{center}
\end{figure}

In this scenario, the acceptable difference overlaps with the original 95\% prediction intervals for understatements and overstatements, and the actual level of assurance obtained is now reduced.

\bigskip
It should be noted that the auditor determines the acceptable levels of risks incurred. The desired level of assurance is related to the combination of inherent risk (base, elevated or significant) and control risk, and may range from 95\% for a significant inherent risk without reliance on controls to much lower levels for a base inherent risk with control reliance. If the auditor wants to primarily control the efficiency risk, the risk can be limited to, for example, 1\% or 5\%. However, the choice of the method determines \emph{which} risk is controlled. If the auditor decides to control efficiency risk, then the desired level of assurance cannot be controlled simultaneously; it is a function of the precision of the expectation. Similarly, if the auditor decides to control the effectiveness risk by setting the desired level of assurance, the efficiency risk is a function of the expectation precision. 
% As remarked before, the precision of a predictive method cannot be easily controlled. 

\bigskip
We conclude the discussion on the acceptable difference with six arguments, five of which support the primary control of efficiency risk.

\begin{enumerate}
	\item Considering materiality.
	\item Control of precision.
	\item Cumulative audit evidence.
	\item The effect of inefficiency.
	\item Practical examples from professional literature.
	\item Retroactive adjustment of the desired level of assurance.
\end{enumerate}

\paragraph{Considering materiality} The Standards mention that the determination of the acceptable difference is influenced by materiality and the desired level of assurance. In the previous section it was demonstrated that the precision of the expectation should also be considered when making this determination. Because the standards do not require precision to be calculated, they are mostly silent about the way the acceptable difference is determined, or how the level of assurance is related to the precision of the expectation. For example, AAG-DATA 3.59 states that 'as with other substantive analytical procedures, professional judgment is used to determine what is considered a significant difference between the auditor's prediction and the recorded amount. Performance materiality is considered when making this judgment.' 

One exception is in AAG-DATA B.23, where it is specifically mentioned that the desired level of assurance can be reduced if the precision of the model is poor. 'In evaluating the level of assurance that a regression analysis may provide, one consideration is how the standard error obtained from the regression compares with the auditor's performance materiality. In this example, if the standard error level obtained is less than performance materiality, this provides further confidence regarding the use of regression. However, if the standard error is a high percentage of performance materiality, the auditor would consider limiting how much assurance the auditor intends to derive from the regression. Regression analyses that factor in performance materiality and the auditor's desired level of assurance would ordinarily adjust achieved assurance for the standard error.'

Appendix B of the AAG-DATA provides a fully worked-out example of regression analysis. In paragraphs B.38-39, it is suggested that the prediction interval be used to determine the range in which recorded revenues are expected to fall. In other words, the example determines the acceptable difference by controlling for efficiency risk. The role of materiality in this determination has not been discussed.

Clearly, determining the acceptable difference based on materiality and the desired level of assurance is in line with the auditing standards AU-C 520, ISA 520, and AS 2305. This is exactly how the minimum sample size of a substantive test of details is determined, and it seems as if the standard setters looked for inspiration in ISA / AU-C 530 when drafting ISA / AU-C 520. The method is straightforward for determining the level of assurance obtained from the SAP. However, the consequences of choosing this method have not been discussed. As we have seen in the discussion of Figure \ref{fig:method_detection_less_precise}, a mismatch between the desired level of assurance and the precision of the expectation can lead to a situation where every single difference, or a large portion of them, requires follow-up, thereby increasing efficiency risk to practically unacceptable levels.

\paragraph{Control of precision} Professional literature offers a solution for cases where the model yields expectations that are not precise: the auditor can improve the precision of the expectation by disaggregating the data, adding predictors\footnote{A \textit{predictor} is an independent variable used in the explanation of variance of the dependent variable.}, or increasing the number of observations with which the model is estimated. However, this may be challenging in practice. Disaggregated data may not always be available for the dependent variable and all predictors, the number of predictors used is limited by the number of historical observations, and historical information may not be easily available. As a result, it is much more difficult to improve the precision of the expectation in an SAP than it is, for example, to reduce the allowance for sampling risk in a statistical sample. Increasing the sample size is an easy solution that reduces the allowance for sampling risk.

\paragraph{Cumulative audit evidence}The auditor's\footnote{AU-C 520.A7} substantive procedures to address the assessed risk of material misstatement for relevant assertions may be ToDs, SAPs, or a combination of both. It does not make sense to express the desired level of assurance for each substantive procedure that is executed. When multiple substantive procedures are planned, every single one of them provides some level of assurance. It is therefore logical that the desired level of assurance is only expressed for the last procedure executed, and that for procedures executed earlier, the level of assurance \emph{obtained} is considered (in combination with the assessed inherent risk and control reliance) as the target for the last procedure. Typically, SAPs are executed before the ToDs. Following this line of reasoning, it makes sense to determine the acceptable difference based on the precision of the expectation.

\paragraph{The effect of inefficiency} A large occurrence of false positives may have an adverse effect on auditors' willingness to apply SAPs. For example, if further investigation is often required, but the account is not misstated, the auditor may become frustrated and lose confidence in the procedure, reverting to a single ToD.

\paragraph{Practical examples from professional literature} The examples in the AICPA practice guides determine the acceptable difference using only the expectation precision. These examples do not refer to materiality in this determination. For example, The AICPA Audit Guide to Audit Data Analytics provides an example of the application of regression analysis, in which the acceptable difference is explained in paragraph B.40, using an efficiency risk of 5\%. In the \textit{Case: On the Go Stores}, the analysis point to further investigate the stores that exhibit the largest positive residuals.\footnote{AAG-ANP A.67 and A.69} 

\paragraph{Retroactive adjustment of the desired level of assurance} Another question arises when the precision of the expectation is not aligned with the desired level of assurance. When the auditor notices that many differences require follow-up, there may be a risk that the auditor retroactively adjusts the desired level of assurance to avoid additional work. This phenomenon is referred to as the \textit{anchoring bias}. While we advocate that the precision of the expectation is taken into consideration when determining the acceptable difference, a procedure that adjusts the acceptable difference (thereby reducing the level of assurance obtained) based on the actual differences should be frowned upon.

% subsection acceptable_difference (end)

\subsection{Expectation development} % (fold)
\label{sub:expectation_development}

Once we have satisfied ourselves that the conditions of suitability of the SAP and reliability of the data have been met, the standards require us\footnote{ISA / AU-C 520.5c} to develop an expectation, and evaluate the precision of this expectation.

The semantics used in the Standards and other professional literature may be confusing for data scientists. 
For an auditor, \hlblue{expectations}\index{expectation} are predictions of recorded accounts or ratios, whereas, for a data scientist, the \hlblue{expected value}\index{expected value}\index{value!expected} is the mean of the sampling distribution of a variable $X$, generally denoted as $E(X)$. 
\hlblue{Prediction}\index{prediction} is commonly used in the sense of either using data to make a statement about the future, as in forecasting, or (closer to our application of analytical procedures) using data to 'predict' outcomes that have already occurred. 
In such cases, prediction involves statistical associations among past values of variables. 
For an auditor, prediction and expectation can be used interchangeably. 
A data scientist identifies two concepts.

Expectations are developed by identifying plausible relationships (e.g., store square footage and retail sales) that are reasonably expected to exist, based on our understanding of the client and the industry in which the client operates. We can select from a variety of data sources to create expectations. For example, we may use prior period information, management budgets or forecasts, industry data, or nonfinancial data.

In some cases, the development of expectations can be simple. Take, for instance, the account ‘Share Capital.’ With our understanding of the client a comparison of this year’s amount with the previous year’s amount could be sufficient. At the back of our mind is a model that says: ‘if I have no knowledge or information that my client issued additional shares during the current year, my expectation is that the amount of Share Capital is the same as last year.’ A similar argument can be used for relatively stable accounts, such as Rent: we would expect that the monthly amount is constant over time,  except for an annual increase.

% subsection expectation_development (end)

% section application_of_the_sap (end)

\section{Evaluating and responding to the SAP} % (fold)
\label{sec:evaluation_and_response}

We now look at the requirements and guidance concerning observations that are not acceptable. These are 'fluctuations or relationships that, on the one hand, are inconsistent with other relevant information, or on the other hand, differ from expected values by a significant amount'.\footnote{ISA / AU-C 520.7} For very precise analyses, the acceptable range is very narrow, so it is possible that the difference observed is outside the range, but not material. We still need to follow up this difference because it questions the model used to form expectations.

We investigate differences by making inquiries of management or by performing other audit procedures. Audit evidence relevant to management's responses may be obtained by evaluating those responses, taking into account our understanding of the entity and its environment and with other audit evidence obtained during the course of the audit. We should ordinarily corroborate explanations for the significant differences by obtaining sufficient and appropriate audit evidence. The procedures used to corroborate the explanation depend on the nature of the explanation, nature of the account balance, and results of other substantive procedures. To corroborate this explanation, one or more of the following techniques may be used:
\begin{compactitem}
\item Inquiries of persons outside the client's organization. For example, we may confirm the discounts received with major suppliers or agree to changes in commodity prices with a commodity exchange or the financial press.
\item Inquiries of independent persons inside the client's organization. For example, an explanation received from the financial controller for an increase in advertising expenditure may be corroborated by the marketing director. It is normally inappropriate to corroborate ex\-pla\-nations only through discussions with other accounting personnel.
\item Evidence obtained from other auditing procedures. Sometimes, the results of other auditing procedures (particularly those performed on the data used to develop expectations) are sufficient to corroborate an explanation.
\item Examination of supporting evidence. The auditor may examine supporting documentary evidence of transactions to corroborate the explanations. For example, if an increase in the cost of sales in one month is attributed to an unusually large sales contract, the auditor might examine supporting documentation such as sales contracts and delivery dockets.
\end{compactitem}

\bigskip
It is usually impractical to identify the factors that explain the exact amount of a difference identified for investigation. However, we should attempt to quantify the portion of the difference for which plausible explanations can be obtained and, where appropriate, corroborated and determine the amount that cannot be explained is sufficiently small (typically: back inside the acceptable difference after adjusting for the explained amount), so we can conclude on the absence of material misstatement.

If a reasonable explanation cannot be obtained, we should determine whether uncorrected misstatements are material, individually or in the aggregate. In making this determination, we should consider 

\begin{compactitem}
\item the size and nature of the misstatements, both in relation to particular classes of transactions, account balances, or disclosures and the financial statements as a whole, and the particular circumstances of their occurrence; and 
\item the effect of uncorrected misstatements related to prior periods on the relevant classes of transactions, account balances, or disclosures and the financial statements as a whole.
\end{compactitem}

% section evaluation_and_response (end)

\section{Documentation} % (fold)
\label{sec:documentation}

There is one last difference between the US Standard and the ISA. AU-C 520.8 provides an additional requirement for the documentation of the procedure. In particular, we should document 
\begin{compactitem}
\item how we formed the expectation;
\item the comparison to the recorded amount or amounts;
\item the calculation of the difference(s); and 
\item the audit procedures performed for those differences that were outside the acceptable difference. 
\end{compactitem}

The fact that the ISA does not specifically state these documentation requirements does not imply that we should not document our procedures. 

% section documentation (end)

% chapter analytical_procedures (end)